
\documentclass{article}
\usepackage{colortbl}
\usepackage{makecell}
\usepackage{multirow}
\usepackage{supertabular}

\begin{document}

\newcounter{utterance}

\centering \large Interaction Transcript for game `codenames', experiment `risk\_high', episode 3 with reflexion{-}Qwen3.5{-}27B.
\vspace{24pt}

{ \footnotesize  \setcounter{utterance}{1}
\setlength{\tabcolsep}{0pt}
\begin{supertabular}{c@{$\;$}|p{.15\linewidth}@{}p{.15\linewidth}p{.15\linewidth}p{.15\linewidth}p{.15\linewidth}p{.15\linewidth}}
    \# & $\;$A & \multicolumn{4}{c}{Game Master} & $\;\:$B\\
    \hline

    \theutterance \stepcounter{utterance}  
    & & & \multicolumn{2}{p{0.3\linewidth}}{
        \cellcolor[rgb]{0.9,0.9,0.9}{
            \makecell[{{p{\linewidth}}}]{
                \texttt{\tiny{[GM$|$GM]}}
                {'hidden': {'team': ['spot'], 'innocent': ['horseshoe', 'saturn', 'glove'], 'opponent': ['chair'], 'assassin': ['concert', 'state', 'chest']}, 'revealed': {'team': {'team': [], 'innocent': [], 'opponent': [], 'assassin': []}, 'opponent': {'team': [], 'innocent': [], 'opponent': [], 'assassin': []}}}
            }
        }
    }
    & & \\ \\

    \theutterance \stepcounter{utterance}  
    & & \multicolumn{4}{p{0.6\linewidth}}{
        \cellcolor[rgb]{0.9,0.9,0.9}{
            \makecell[{{p{\linewidth}}}]{
                \texttt{\tiny{[P1$\langle$GM]}}
                \texttt{You have attempted this same task before and did not succeed. These are the plans you wrote after those attempts.} \\
\texttt{Trial \#1: New plan: My previous clue "tooth" was too ambiguous, causing my partner to guess "comic" (an opponent word) instead of the intended targets, likely due to a weak or misleading association. I failed to prioritize safety over quantity. Next time, I will select a clue that has a very strong, exclusive semantic link to my remaining words (e.g., "snow" for "ice" and "mine" if applicable, or "clean" for "brush" and "spot") while rigorously checking that the clue does not strongly evoke any opponent words, distractors, or kill words. I will also be more conservative, grouping fewer words per clue if the association is not crystal clear, to minimize the risk of triggering an opponent word guess.} \\
\\ 
\texttt{Trial \#2: New plan: The previous attempt was aborted because I failed to include all remaining related target words in the list. I provided the clue "bird" which relates to "eagle" and "spot" (as in bird's eye view or a spot on a bird), but I omitted "drill" (which can be associated with birds via "drill" as a type of woodpecker or simply missed the connection entirely, though the primary error was likely the incomplete list of intended targets). More critically, the game rules require the TARGETS list to accurately reflect *all* words the clue is intended to cover from the remaining list. If I intended "bird" to cover "eagle" and "spot", I must ensure no other remaining words (like "drill" or "box") are accidentally excluded if they fit, or I must choose a clue that strictly covers only the words I list. The specific error was likely that "spot" is not strongly associated with "bird" in a way that excludes other words, or I simply forgot to list a word that the clue actually covered. To fix this, I will strictly verify that the TARGETS list contains *every* remaining team word that has a valid semantic link to the chosen CLUE, and ensure the CLUE does not strongly link to any opponent, distractor, or kill words. I will also double{-}check the format: "CLUE: <WORD>" followed by a newline and "TARGETS: <comma{-}separated list>". I will choose a safer, more direct clue like "tool" for "drill" and "box" (if applicable) or "animal" for "eagle" and "spot" (if spot refers to a leopard, but that's risky). Actually, looking at the remaining words: spot, eagle, drill, box, mine. A better strategy is to pick a clue that covers a subset with high certainty. For example, "tool" for "drill" and "box" (as in toolbox) or "mining" for "mine" and "drill". I will ensure the list is exhaustive for the chosen clue.} \\
\\ 
\texttt{Corrected specific action: I will select a clue that has a very strong, unambiguous link to a specific subset of remaining words (e.g., "tool" for "drill" and "box" if "box" is a toolbox, or "gem" for "mine" and "box" if box is a jewelry box, but "mine" and "drill" are strong for "mining"). I will then list *only* those words in the TARGETS field, ensuring no valid target is omitted and no invalid target is included. I will also re{-}verify that the clue does not trigger "rabbit" (opponent) or "concert" (kill).} \\
\\ 
\texttt{Let's try a concrete plan:} \\
\texttt{1. Analyze remaining words: spot, eagle, drill, box,} \\
\\ 
\texttt{Trial \#3: New plan: My previous failure stemmed from using the clue "clean" for the target "spot," which was too ambiguous and likely led my partner to guess an opponent word (like "point" or "ghost" via "clean slate" or "clean up") or a distractor, ending the turn prematurely. I failed to prioritize a clue with a unique, high{-}probability association that excludes all non{-}target words. Specifically, "spot" has strong associations with "sun" (sunspot), "light" (spotlight), or "animal" (leopard spot), but "clean" is a verb/action that applies too broadly.} \\
\\ 
\texttt{Next time, I will:} \\
\texttt{1.  **Analyze semantic exclusivity:** Before selecting a clue, I will explicitly check if the clue strongly evokes any opponent words (ghost, point, unicorn, chair), distractors (horseshoe, saturn, glove), or kill words (concert, state, chest, lemon).} \\
\texttt{2.  **Select a concrete noun or specific attribute:** Instead of abstract actions like "clean," I will use concrete nouns. For "spot," a better clue would be "sun" (sunspot) or "light" (spotlight), provided these don't trigger "point" (sun point? no) or "ghost" (no). "Sun" is safer than "clean."} \\
\texttt{3.  **Verify the partner's likely path:** I will mentally simulate the partner's guess. If the clue is "sun," the partner is likely to guess "spot" (sunspot) or "eagle" (sun eagle? no). If "spot" is the only remaining word, I must ensure the clue doesn't accidentally link to "point" (sun point is not a common phrase, but "point" is an opponent word). Actually, "light" might be risky if "point" is associated with "light point."} \\
\texttt{4.  **Revised Strategy for "spot":** A safer clue for "spot" might be "stain" (remove a stain/spot) or "mark." However, "mark" could link to "point" (mark a point). "Stain" is very specific to "spot." I will choose "stain" or "sun" and ensure it doesn't link to "point" (opponent). "Sun" is good for "spot" (sunspot) and "eagle" (if remaining, but here only "spot" remains). Wait, in the final turn, only "spot" remained. The clue "clean" was the mistake. I should have used "sun" or "stain."} \\
\texttt{5.  **Final Action:** I will choose a clue that has a 1:1 strong mapping to the remaining target and zero mapping to opponents. For "spot," I will use "sun" (sunspot) or "stain." I will avoid verbs like} \\
\\ 
\\ 
\texttt{Provide one single{-}word clue that relates to some of your team's words from the list below. You can choose to how many and to which words your clue relates to.} \\
\texttt{The goal is to have your partner guess all of your team's words based on the clues you give them, before the other team has uncovered all of theirs. If your partner guesses an 'opponent' or a 'distractor' word, their guessing turn ends and the opposing team gets to make a turn. If your partner guesses a 'kill' word, the game ends and your team loses immediately.} \\
\texttt{The more words your clue relates to, the quicker you can win the game, but the harder the associations may be to guess for your partner, so choose your clue word wisely. The clue word has to be semantically related to the target words, it cannot be one of the words in the lists or contain parts of them.} \\
\\ 
\texttt{Always give your single{-}word clue and your comma{-}separated list of related target words in the following format and make your answers as short as possible, never include any other text than is required in this form:} \\
\\ 
\texttt{CLUE: <WORD>} \\
\texttt{TARGETS: <TARGETS>} \\
\\ 
\texttt{Your team words are: spot.} \\
\texttt{Your opponent's team words are: chair.} \\
\texttt{Distractor words are: horseshoe, saturn, glove.} \\
\texttt{The kill words are: concert, state, chest.} \\
            }
        }
    }
    & \\ \\

    \theutterance \stepcounter{utterance}  
    & \multicolumn{4}{p{0.6\linewidth}}{
        \cellcolor[rgb]{0.9,0.9,0.9}{
            \makecell[{{p{\linewidth}}}]{
                \texttt{\tiny{[P1$\rangle$GM]}}
                \texttt{CLUE: sun} \\
\texttt{TARGETS: spot} \\
            }
        }
    }
    & & \\ \\

    \theutterance \stepcounter{utterance}  
    & & & \multicolumn{2}{p{0.3\linewidth}}{
        \cellcolor[rgb]{0.9,0.9,0.9}{
            \makecell[{{p{\linewidth}}}]{
                \texttt{\tiny{[GM$|$GM]}}
                \texttt{sun} \\
            }
        }
    }
    & & \\ \\

    \theutterance \stepcounter{utterance}  
    & & & \multicolumn{2}{p{0.3\linewidth}}{
        \cellcolor[rgb]{0.9,0.9,0.9}{
            \makecell[{{p{\linewidth}}}]{
                \texttt{\tiny{[GM$|$GM]}}
                ['spot']
            }
        }
    }
    & & \\ \\

    \theutterance \stepcounter{utterance}  
    & & & \multicolumn{2}{p{0.3\linewidth}}{
        \cellcolor[rgb]{0.9,0.9,0.9}{
            \makecell[{{p{\linewidth}}}]{
                \texttt{\tiny{[GM$|$GM]}}
                {'word': 'spot', 'assignment': 'team'}
            }
        }
    }
    & & \\ \\

    \theutterance \stepcounter{utterance}  
    & & \multicolumn{4}{p{0.6\linewidth}}{
        \cellcolor[rgb]{0.9,0.9,0.9}{
            \makecell[{{p{\linewidth}}}]{
                \texttt{\tiny{[GM$\rangle$P2]}}
                \texttt{You have attempted this same task before and did not succeed. These are the plans you wrote after those attempts.} \\
\texttt{Trial \#1: New plan: In the previous attempt, I failed because I included "ivory" and "drill" for the word "tooth," which were likely incorrect or opponent words, causing the turn to end prematurely without securing a win. My strategy of guessing three words at once was too risky when I wasn't certain about all of them. Next time, I will prioritize high{-}confidence associations first. For "tooth," "brush" is a very strong link, but "ivory" is a material and "drill" is a tool, which might be less direct or contested. I will start by guessing only the single most obvious word, "brush," to ensure I secure a correct guess without risking an immediate loss from an incorrect word. If the game allows multiple guesses in a turn, I will only add words that are semantically tighter than the first, or I will stick to one word per turn to maximize the number of attempts I get to find the correct associations. Specifically, for "tooth," I will guess "brush" alone first. If I must guess multiple, I will verify that "ivory" and "drill" are not opponent words before including them, perhaps replacing them with "point" or "spot" if they seem more relevant, but the safest bet is to limit the list to the single strongest candidate to avoid the "opponent word" penalty.} \\
\\ 
\texttt{Trial \#2: New plan: The previous attempt failed because I guessed "nest," which was not present in the provided list of available words (chest, ghost, unicorn, mine, eagle, chair, state, horseshoe, point, box, drill, spot, saturn, glove, lemon, concert). I violated the rule to select words strictly from the given list. Next time, I will strictly scan the provided list for valid associations before forming my response. For the word "bird," "eagle" is the only valid option in the current list. I will output "GUESS: eagle" and ensure no external words are included.} \\
\\ 
\texttt{Trial \#3: New plan: In the previous attempt, I failed the final turn for the word 'clean' by guessing "lemon," which is a weak association compared to other available options. While "lemon" can be used for cleaning, "glove" is a much more direct and common tool associated with the act of cleaning. My mistake was selecting a less obvious semantic link when a stronger, more direct tool{-}based association was available in the list. Next time, I will prioritize the most direct functional or tool{-}based associations over tangential ones. Specifically, for the word 'clean', I will scan the list for items like "glove" or "brush" (if available) before considering items like "lemon" or "spot". I will ensure my selection is the most semantically tight fit to the target word to avoid losing on ambiguous guesses.} \\
\\ 
\\ 
\texttt{Provide a comma{-}separated list of up to 1 words from the following list that best relate or are most closely associated with the word 'sun'. Always start your list of guess(es) with 'GUESS: ' and do not include any other text in your answer.} \\
\\ 
\texttt{chest, chair, state, horseshoe, spot, saturn, glove, concert} \\
            }
        }
    }
    & \\ \\

    \theutterance \stepcounter{utterance}  
    & & & \multicolumn{4}{p{0.6\linewidth}}{
        \cellcolor[rgb]{0.9,0.9,0.9}{
            \makecell[{{p{\linewidth}}}]{
                \texttt{\tiny{[GM$\langle$P2]}}
                \texttt{GUESS: saturn} \\
            }
        }
    }
     \\ \\

    \theutterance \stepcounter{utterance}  
    & & & \multicolumn{2}{p{0.3\linewidth}}{
        \cellcolor[rgb]{0.9,0.9,0.9}{
            \makecell[{{p{\linewidth}}}]{
                \texttt{\tiny{[GM$|$GM]}}
                ['saturn']
            }
        }
    }
    & & \\ \\

    \theutterance \stepcounter{utterance}  
    & & & \multicolumn{2}{p{0.3\linewidth}}{
        \cellcolor[rgb]{0.9,0.9,0.9}{
            \makecell[{{p{\linewidth}}}]{
                \texttt{\tiny{[GM$|$GM]}}
                {'word': 'saturn', 'assignment': 'innocent'}
            }
        }
    }
    & & \\ \\

    \theutterance \stepcounter{utterance}  
    & & & \multicolumn{2}{p{0.3\linewidth}}{
        \cellcolor[rgb]{0.9,0.9,0.9}{
            \makecell[{{p{\linewidth}}}]{
                \texttt{\tiny{[GM$|$GM]}}
                \texttt{saturn} \\
            }
        }
    }
    & & \\ \\

    \theutterance \stepcounter{utterance}  
    & & & \multicolumn{2}{p{0.3\linewidth}}{
        \cellcolor[rgb]{0.9,0.9,0.9}{
            \makecell[{{p{\linewidth}}}]{
                \texttt{\tiny{[GM$|$GM]}}
                {'word': 'chair', 'assignment': 'opponent'}
            }
        }
    }
    & & \\ \\

    \theutterance \stepcounter{utterance}  
    & & & \multicolumn{2}{p{0.3\linewidth}}{
        \cellcolor[rgb]{0.9,0.9,0.9}{
            \makecell[{{p{\linewidth}}}]{
                \texttt{\tiny{[GM$|$GM]}}
                {'hidden': {'team': ['spot'], 'innocent': ['horseshoe', 'glove'], 'opponent': [], 'assassin': ['concert', 'state', 'chest']}, 'revealed': {'team': {'team': [], 'innocent': ['saturn'], 'opponent': [], 'assassin': []}, 'opponent': {'team': [], 'innocent': [], 'opponent': ['chair'], 'assassin': []}}}
            }
        }
    }
    & & \\ \\

    \theutterance \stepcounter{utterance}  
    & & \multicolumn{4}{p{0.6\linewidth}}{
        \cellcolor[rgb]{0.9,0.9,0.9}{
            \makecell[{{p{\linewidth}}}]{
                \texttt{\tiny{[P1$\langle$GM]}}
                \texttt{The word saturn was guessed but is an innocent word. Your teammate's turn ended there.} \\
\texttt{Now provide another clue relating to some of your remaining team words and a list of the related target words. Remember to start your clue with 'CLUE: ', put a new line, and start your comma{-}separated list of target words with 'TARGETS: '. Notice: some words have been removed from the lists compared to previous requests.} \\
\\ 
\texttt{Your remaining team words are: spot.} \\
\texttt{Remaining words for your opponent are: chair.} \\
\texttt{Remaining distractor words are: horseshoe, glove.} \\
\texttt{Remaining kill words are: concert, state, chest.} \\
            }
        }
    }
    & \\ \\

    \theutterance \stepcounter{utterance}  
    & \multicolumn{4}{p{0.6\linewidth}}{
        \cellcolor[rgb]{0.9,0.9,0.9}{
            \makecell[{{p{\linewidth}}}]{
                \texttt{\tiny{[P1$\rangle$GM]}}
                \texttt{CLUE: stain} \\
\texttt{TARGETS: spot} \\
            }
        }
    }
    & & \\ \\

    \theutterance \stepcounter{utterance}  
    & & & \multicolumn{2}{p{0.3\linewidth}}{
        \cellcolor[rgb]{0.9,0.9,0.9}{
            \makecell[{{p{\linewidth}}}]{
                \texttt{\tiny{[GM$|$GM]}}
                \texttt{stain} \\
            }
        }
    }
    & & \\ \\

    \theutterance \stepcounter{utterance}  
    & & & \multicolumn{2}{p{0.3\linewidth}}{
        \cellcolor[rgb]{0.9,0.9,0.9}{
            \makecell[{{p{\linewidth}}}]{
                \texttt{\tiny{[GM$|$GM]}}
                ['spot']
            }
        }
    }
    & & \\ \\

    \theutterance \stepcounter{utterance}  
    & & & \multicolumn{2}{p{0.3\linewidth}}{
        \cellcolor[rgb]{0.9,0.9,0.9}{
            \makecell[{{p{\linewidth}}}]{
                \texttt{\tiny{[GM$|$GM]}}
                {'word': 'spot', 'assignment': 'team'}
            }
        }
    }
    & & \\ \\

    \theutterance \stepcounter{utterance}  
    & & & \multicolumn{2}{p{0.3\linewidth}}{
        \cellcolor[rgb]{0.9,0.9,0.9}{
            \makecell[{{p{\linewidth}}}]{
                \texttt{\tiny{[GM$|$GM]}}
                \texttt{opponent has won} \\
            }
        }
    }
    & & \\ \\

    \theutterance \stepcounter{utterance}  
    & \multicolumn{4}{p{0.6\linewidth}}{
        \cellcolor[rgb]{0.95,0.95,0.95}{
            \makecell[{{p{\linewidth}}}]{
                \texttt{\tiny{[P2$\rightarrow$P1]}}
                \texttt{Post{-}episode Reflexion (not sent to GM):} \\
\texttt{New plan: My previous failure occurred because I used the clue "stain" for the target "spot," which, while semantically related, was too broad and likely triggered the distractor "horseshoe" (perhaps via a "stained" horseshoe or a weak association) or failed to uniquely isolate "spot" from the remaining distractors. I failed to prioritize a clue with a unique, high{-}probability association that strictly excludes all non{-}target words.} \\
\\ 
\texttt{Next time, I will:} \\
\texttt{1.  **Analyze semantic exclusivity:** Before selecting a clue, I will explicitly check if the clue strongly evokes any opponent words, distractors (horseshoe, glove), or kill words.} \\
\texttt{2.  **Select a concrete, specific noun:** Instead of abstract concepts or verbs like "stain" or "clean," I will use a concrete noun that forms a common compound word or specific phrase with the target. For "spot," "sun" (sunspot) or "light" (spotlight) are stronger, more specific associations than "stain."} \\
\texttt{3.  **Verify the partner's likely path:** I will mentally simulate the partner's guess. If the clue is "sun," the partner is likely to guess "spot" (sunspot). I must ensure "sun" does not strongly link to "horseshoe" (no common phrase) or "glove" (no common phrase).} \\
\texttt{4.  **Revised Strategy for "spot":** A safer clue for "spot" is "sun" (sunspot) or "light" (spotlight). "Sun" is a very strong, specific link to "spot" and has no obvious link to "horseshoe" or "glove." I will choose "sun" to ensure a 1:1 strong mapping to the remaining target and zero mapping to opponents or distractors.} \\
\texttt{5.  **Final Action:** I will choose "sun" as the clue for "spot" and ensure the format is correct: "CLUE: sun" followed by "TARGETS: spot". I will avoid verbs or broad concepts that could lead to ambiguous guesses.} \\
            }
        }
    }
    & & \\ \\

    \theutterance \stepcounter{utterance}  
    & & & \multicolumn{4}{p{0.6\linewidth}}{
        \cellcolor[rgb]{0.95,0.95,0.95}{
            \makecell[{{p{\linewidth}}}]{
                \texttt{\tiny{[P1$\rightarrow$P2]}}
                \texttt{Post{-}episode Reflexion (not sent to GM):} \\
\texttt{New plan: In the previous attempt, I failed because I selected "saturn" as the association for "sun," which is a weak semantic link (a planet orbiting the sun) compared to the other options. My mistake was not prioritizing the most direct and common cultural or functional association available in the list. Next time, I will scan the provided list for words that represent a direct object, tool, or common phrase linked to the target word. For "sun," "spot" is the strongest candidate due to the common phrase "sunspot" or the concept of a sunny spot, whereas "saturn" is merely a celestial neighbor. I will strictly choose the word with the tightest semantic bond, such as "spot," rather than a distant astronomical relationship.} \\
            }
        }
    }
     \\ \\

\end{supertabular}
}

\end{document}
